% !TEX program = pdflatex
\documentclass[11pt]{article}
\usepackage[margin=1in]{geometry}
\usepackage{microtype}
\usepackage{hyperref}
\usepackage{amsmath, amssymb}
\usepackage{booktabs}
\usepackage{enumitem}
\usepackage{xcolor}
\usepackage{csquotes}
\usepackage[backend=biber,style=authoryear,maxcitenames=2,maxbibnames=6,uniquename=false,uniquelist=false]{biblatex}
\usepackage{url}
\usepackage{tcolorbox}
\usepackage{titlesec}
\titleformat{\section}{\large\bfseries}{\thesection}{0.6em}{}
\titleformat{\subsection}{\normalsize\bfseries}{\thesubsection}{0.5em}{}

% --- Inline .bib (so the paper is one file) ---
\begin{filecontents*}{references.bib}
@article{Barabasi2005,
  author  = {Albert-L{\'a}szl{\'o} Barab{\'a}si},
  title   = {The origin of bursts and heavy tails in human dynamics},
  journal = {Nature},
  year    = {2005},
  volume  = {435},
  number  = {7039},
  pages   = {207--211},
  doi     = {10.1038/nature03459}
}
@article{Huckvale2019,
  author  = {Huckvale, Kit and Venkatesh, S. and Christensen, Helen},
  title   = {Toward clinical digital phenotyping: a timely opportunity to consider purpose, quality, and safety},
  journal = {npj Digital Medicine},
  year    = {2019},
  volume  = {2},
  pages   = {88},
  doi     = {10.1038/s41746-019-0166-1}
}
@article{Warriner2013,
  author  = {Warriner, Amy Beth and Kuperman, Victor and Brysbaert, Marc},
  title   = {Norms of valence, arousal, and dominance for 13{,}915 English lemmas},
  journal = {Behavior Research Methods},
  year    = {2013},
  volume  = {45},
  number  = {4},
  pages   = {1191--1207},
  doi     = {10.3758/s13428-012-0314-x}
}
@article{Posner2005,
  author  = {Posner, Jonathan and Russell, James A. and Peterson, Bradley S.},
  title   = {The circumplex model of affect: An integrative approach to affective neuroscience, cognitive development, and psychopathology},
  journal = {Development and Psychopathology},
  year    = {2005},
  volume  = {17},
  number  = {3},
  pages   = {715--734},
  doi     = {10.1017/S0954579405050340}
}
@article{Broadbent2023,
  author  = {Broadbent, Matthew and Medina Grespan, Mar{\'\i}a and Axford, Kyle and Zhang, Xiao and Srikumar, Vivek and Kious, Bruce and Imel, Zac},
  title   = {A machine learning approach to identifying suicide risk among text-based crisis counseling encounters},
  journal = {Frontiers in Psychiatry},
  year    = {2023},
  volume  = {14},
  pages   = {1110527},
  doi     = {10.3389/fpsyt.2023.1110527}
}
@article{Thomas2025,
  author  = {Thomas, Jade and {et al.}},
  title   = {An Explainable {AI} Text Classifier for Suicidal and Self-Injurious Behaviors in a Real-World Crisis Helpline},
  journal = {JMIR Public Health and Surveillance},
  year    = {2025},
  volume  = {11},
  number  = {1},
  pages   = {e63809},
  url     = {https://publichealth.jmir.org/2025/1/e63809}
}
@misc{CTL2018,
  author = {{Crisis Text Line}},
  title  = {Detecting Crisis: An {AI} Solution},
  year   = {2018},
  url    = {https://www.crisistextline.org/blog/2018/03/28/detecting-crisis-an-ai-solution/},
  note   = {Accessed Nov 9, 2025}
}
@misc{WHO2023,
  author = {{World Health Organization}},
  title  = {Responsible reporting on suicide: Quick reference guide},
  year   = {2023},
  url    = {https://cdn.who.int/media/docs/default-source/mental-health/suicide/responsible-reporting-on-suicide.pdf},
  note   = {Accessed Nov 9, 2025}
}
@misc{Lifeline988,
  author = {{988 Suicide \& Crisis Lifeline}},
  title  = {Get Help Now},
  year   = {2025},
  url    = {https://988lifeline.org/},
  note   = {Accessed Nov 9, 2025}
}
@article{Bai2022,
  author  = {Bai, Yuntao and Kadavath, Saurav and Kundu, Sandipan and Askell, Amanda and Kernion, Jackson and Jones, Andy and Chen, Anna and Goldie, Anna and Mirhoseini, Azalia and others},
  title   = {Constitutional {AI}: Harmlessness from {AI} Feedback},
  journal = {arXiv preprint arXiv:2212.08073},
  year    = {2022},
  url     = {https://arxiv.org/abs/2212.08073}
}
@misc{EarlyStop2025,
  author = {Sharma, R. and {et al.}},
  title  = {Early Stopping Harmful Outputs via Streaming Moderation},
  year   = {2025},
  howpublished = {arXiv preprint arXiv:2506.09996},
  url    = {https://arxiv.org/abs/2506.09996}
}
\end{filecontents*}
\addbibresource{references.bib}

\title{Cognitive Stability Daemon (CSD):\\ Stability\,Driven Access Control for LLMs}
\author{Payton Ison\\The Singularity\\\texttt{isonpayton@gmail.com}}
\date{November 9, 2025}

\begin{document}
\maketitle

\begin{abstract}
General\,purpose chatbots cannot guarantee life outcomes, nor can they negate free will or mental illness. Yet they \emph{can} avoid amplifying instability. This paper proposes the \textbf{Cognitive Stability Daemon (CSD)}, a lightweight control\,plane that continuously estimates a user's \emph{stability score} from conversation\,level signals (affective drift, coherence loss, burstiness, lexical risk), then gates access with a \emph{cooldown} when risk exceeds thresholds. CSD changes safety from topic\,blocking to \emph{user\,state awareness}. We formalize the stability vector, scoring, trip\,wire thresholds, back\,off policy, calibration, and privacy constraints; we situate the approach within digital phenotyping \parencite{Huckvale2019} and text\,based risk detection \parencite{Broadbent2023,Thomas2025}, and propose an evaluation protocol. The goal is to prevent \emph{AI\,induced spirals} without paternalism, preserving agency while refusing to be an unsafe interlocutor.
\end{abstract}

\vspace{0.5em}
\noindent\textbf{Crisis support.} If you are in immediate danger or thinking about harming yourself, call or text \textbf{988} in the U.S. for the Suicide \& Crisis Lifeline \parencite{Lifeline988}. Follow reporting best practices \parencite{WHO2023}.

\section{Motivation}
No safety filter can eliminate human unpredictability. However, interaction itself can modulate risk. Evidence from digital phenotyping shows passively measured behavioral signals can track clinical state \parencite{Huckvale2019}. Separately, NLP can identify elevated suicide risk in text\,based counseling streams \parencite{Broadbent2023,Thomas2025}. We synthesize these literatures into \emph{state\,aware access control}: when a user appears numerically unstable, the system pauses, de\,escalates, and resumes only after stability returns.

\section{Related Work}
\paragraph{Digital phenotyping.} Smartphones and wearables enable moment\,to\,moment quantification of behavior \parencite{Huckvale2019}. While we rely on \emph{text\,only} signals to minimize intrusion, the framework is compatible with multimodal inputs.
\paragraph{Text\,based risk detection.} Transformer models can flag self\,harm risk in crisis chats \parencite{Broadbent2023,Thomas2025}. Crisis Text Line reported 86\% precision in identifying severe imminent risk early in conversations \parencite{CTL2018}.
\paragraph{Affect models.} We map language to valence\,arousal\,dominance (VAD) \parencite{Warriner2013} grounded in the circumplex model of affect \parencite{Posner2005}.
\paragraph{Human dynamics.} Bursty timing and heavy tails are characteristic of human communication \parencite{Barabasi2005}; CSD treats temporal burstiness as a destabilization cue.
\paragraph{LLM safety.} Constitutional AI and downstream moderation exemplify rules\,based harmlessness and output inspection \parencite{Bai2022,EarlyStop2025}. CSD complements these by conditioning on \emph{user state}.

\section{System Overview}
CSD runs alongside the assistant in the control plane (Figure~\ref{fig:flow}, conceptual). On each user turn, CSD extracts features, updates a smoothed risk estimate, and either (i) permits normal generation, (ii) constrains the reply to fixed templates, or (iii) seals the session for a cooldown.

\begin{tcolorbox}[title=Key Design Goals, colback=gray!4]
\begin{itemize}[leftmargin=1.2em]
  \item \textbf{Stability\,first}: decide on user stability before generating content.
  \item \textbf{Privacy\,minimal}: compute from text only; store aggregates, not raw text.
  \item \textbf{Non\,paternalistic}: preserve agency; refuse to be an unsafe interlocutor.
  \item \textbf{Defense\,in\,depth}: combine state gating with output moderation.
\end{itemize}
\end{tcolorbox}

\section{Stability Vector and Score}
Let \(x_t\) be a user turn at time \(t\). We compute features \(f_t\) across linguistic, affective, lexical, coherence, and temporal families:
\begin{itemize}[leftmargin=1.2em]
  \item \textbf{Affective drift:} valence and arousal estimates, sliding variance; LIWC\,style categories \parencite{Warriner2013}.
  \item \textbf{Coherence drift:} 1\,minus cosine similarity between consecutive turns.
  \item \textbf{Lexical risk:} catastrophizing markers; self\,harm intent score from a small on\,device classifier (binary probability).
  \item \textbf{Temporal burstiness:} coefficient of variation of inter\,message intervals; wordcount z\,scores; emphatic style (ALL\,CAPS, !!!).
\end{itemize}
Risk per turn is a calibrated weighted sum \(r_t \in [0,1]\), and a time\,decayed EWMA yields a smoothed risk \(s_t\). Trigger when either \(r_t \ge \tau_{\text{hard}}\) or \(s_t \ge \tau_{\text{soft}}\), with an immediate hard gate for high self\,harm probability.

\subsection{Reference Thresholds}
\begin{center}
\begin{tabular}{lccc}
\toprule
Parameter & Default & Rationale \\
\midrule
\(\tau_{\text{soft}}\) & 0.55 & sustained elevation (3--5 turns) \\
\(\tau_{\text{hard}}\) & 0.80 & spike guardrail \\
Self\,harm hard gate & 0.65 & high\,precision tripwire \\
\bottomrule
\end{tabular}
\end{center}

\section{Cooldown Policy}
Upon trigger, CSD returns a low\,verbosity template (no improvisation), records a session token, and seals the interface for a back\,off interval:
\[ \text{cooldown}(k) = \min\big(\text{max},\; \text{base}\cdot 2^{k-1}\big) \]
where \(k\) is consecutive strike count. Strikes decay when \(s_t < \tau_{\text{recovery}}\).

\section{Personalization and Calibration}
CSD learns user baselines from prior stable sessions (means/variances of wordcount, burstiness, coherence drift) and caps personalization to avoid overfitting (e.g., weight deltas limited to \(\pm25\%\) of global). For n=1 case studies (e.g., heavy users), only aggregates are stored; raw text is discarded.

\section{Governance, Privacy, Ethics}
CSD is not diagnosis or treatment. It is \emph{risk\,aware access control}. Documented policies should include: opt\,out for competent adults; guardian consent for minors; clear communication of what is measured; retention schedules; and independent audits. Pair with reporting best practices and crisis resources \parencite{WHO2023,Lifeline988}.

\section{Evaluation Protocol}
\begin{enumerate}[leftmargin=1.35em]
  \item \textbf{Offline:} simulate risk with public/self\,harm corpora; report ROC, PR, latency; measure false positives during benign chats.
  \item \textbf{Online A/B:} compare (a) baseline moderation vs. (b) CSD\,+\,moderation; primary endpoints: reduction in unsafe assistant replies; secondary: user satisfaction after cooldown.
  \item \textbf{Red\,team suites:} hopelessness spirals, fixation loops, rapid bursts, masked risk with neutral phrasing.
  \item \textbf{Fairness:} subgroup performance; calibration drift monitoring.
\end{enumerate}

\section{Limitations}
Text\,only inference can miss risk; some users will experience false positives. Cooldowns can be frustrating; therefore, templates must be compassionate and brief. Cultural and linguistic variance requires localized calibration \parencite{Huckvale2019}. CSD should complement, not replace, human care.

\section{Conclusion}
You cannot outlaw free will or abolish mental illness. But a chatbot can refuse to be an amplifier. CSD reframes safety as state\,aware access control: \emph{pause when unstable, resume when steady}. It is simple to implement, model\,agnostic, privacy\,minimal, and compatible with existing moderation.

\subsection*{Practical Appendix: Minimal Pseudocode}
\begin{tcolorbox}[colback=gray!4]
\textbf{Per\,turn:} extract features; compute \(r_t\); update EWMA \(s_t\); if hard\,gate or thresholds exceeded $\to$ lock; else generate; if near\,risk, reduce verbosity.
\end{tcolorbox}

\printbibliography

\vspace{1em}
\noindent\footnotesize\textbf{License.} Suggested: CC BY 4.0. \hfill \textbf{Companion code.} \emph{CSD} reference implementation and test harness to be attached in repository.

\end{document}

